\vspace{-5pt}
\section{Related Work}
\vspace{-3pt}
\paragraph{Transition path sampling.}
Traditional approaches to TPS like umbrella sampling \cite{torrie1977nonphysical} and metadynamics \cite{Laio2002a} employ biasing potentials along representative collective variables (CV). However, defining suitable CVs is challenging near transition states, even with automated approaches \cite{Sultan2018AutomatedLearning, Sipka2023}. Meanwhile, shooting techniques \cite{mullen2015easy, borrero2016avoiding, jung2017transition, bolhuis2021transition}, which use a Metropolis-Hastings criterion to sample trajectories, suffer from slow sampling, high rejection rates, and the need for expensive simulations. While ML approaches, including reinforcement learning \cite{das2021reinforcement, rose2021reinforcement, singh2023variational, seong2024transitionpathsamplingimproved, liang2023probing}, differentiable simulations \cite{sipka2023differentiable}, and $h$-transform learning \cite{singh2023variational, du2024doob}, have been used to design CVs or biasing potentials with promising results, they require expensive sampling procedures, must be retrained for each new system of interest, and do not exploit atomistic simulation data. Minimizing the OM action has been used for TPS in low-dimensional systems \cite{faccioli2006dominant, vanden2008geometric, autieri2009dominant, 10.1063/1.3372802, Lee2017}, but has faced computational challenges, such as a lack of integration with modern auto-differentiation frameworks, preventing scalable, gradient-based optimization \cite{a2012dominant}. To our knowledge, our work is the first to propose a gradient-based optimization of the complete OM action, and to connect it to the stochastic dynamics induced by generative models. For a more extensive review of TPS methods, see Appendix \ref{sec:extended_related_work}.

% Meanwhile, data-utilizing methods \cite{Duan2023_oareactdiff, duan2024react, kim2024diffusion}.assume access to a dataset, such as chemical reaction pathways resulting from Nudged Elastic Band (NEB) calculations, and formulate transition state or path prediction as a supervised learning problem. While the presence of data simplifies training and admits the possibility of generalization to unseen systems, the resulting models can only be used for a narrow set of TPS-related tasks, and do not leverage general-purpose representations from other atomistic property prediction tasks, such as energy and force prediction or generative modeling.
\vspace{-7pt}
\paragraph{Atomistic generative models.}
Generative models can produce unbiased, independent samples from the configurational ensemble of molecular systems, pioneered by Boltzmann generators \cite{noe2019boltzmann} and since further developed for proteins \cite{arts2023two, zheng2024predicting, jing2024alphafoldmeetsflowmatching, lewis2024scalable, schreiner2023implicit}, small molecules \cite{huang2024dual, schneuing2024structure,reidenbach2024coarsenconf,igashov2024equivariant}, and materials \cite{zeni2023mattergen, zheng2024predicting, xie2022crystal}. Generative models are typically trained to match the distribution of atomistic configurations from large-scale datasets, including structural databases \cite{bank1971protein} and long-timescale MD simulations \cite{lindorff2011fast, vander2024atlas}. Recent works adapt generative models to produce more diverse samples \cite{corso2023particle}, perform rare event sampling \cite{falkner2023conditioning}, perform MD simulations using the connection between diffusion models and force fields \cite{arts2023two}, and learn generative models directly over trajectories \cite{jing2024generative}.

\vspace{-7pt}
\paragraph{Interpolations in generative models.}
Analagous to TPS, interpolation has been used to evaluate the smoothness and continuity of learned data manifolds and to generate realistic transitions between data points using generative models. While linear interpolation in model latent spaces is known to capture some continuity \cite{kingma2013auto, goodfellow2014generative}, geometric techniques such as geodesic interpolation and optimal transport \cite{arjovsky2017wasserstein,arvanitidis2018latent, lesniak2018distribution, michelis2021linear, struski2023feature, psenka2024representation} better align with the intrinsic manifold structure of the data. Our OM optimization approach can be seen as a novel interpolation mechanism which leverages the inductive bias of stochastic dynamics to generate high-probability transition paths.
